Let $A\subseteq\mathbb{N}$ (here $\mathbb{N}=\{0,1,2,\dots\}$ or $\{1,2,3,\dots\}$; the distinction does not affect ``all sufficiently large integers'').
Say $A$ is an \emph{asymptotic additive basis of order $2$} if there exists $n_0$ such that every integer $n\ge n_0$ can be written as $n=a+a'$ with $a,a'\in A$.
Let
\[
 r_A(n):=(1_A*1_A)(n)=\bigl|\{(a,a')\in A\times A: a+a'=n\}\bigr|,
\]
(the ordered representation function).
A subset $B\subseteq A$ is a \emph{minimal asymptotic basis of order $2$} if $B$ is an asymptotic basis of order $2$ and for every $b\in B$, the set $(B\setminus\{b\})+(B\setminus\{b\})$ fails to contain infinitely many integers.

Question: If $A$ is an asymptotic basis of order $2$ and $r_A(n)\to\infty$ as $n\to\infty$, must $A$ contain a minimal asymptotic basis of order $2$? What if $r_A(n)\ge \varepsilon\log n$ for all large $n$?
